\documentclass[runningheads]{llncs}

\usepackage{graphicx}
\usepackage{amsmath,amssymb}
\usepackage{booktabs}
\usepackage{url}
\usepackage{hyperref}
\usepackage{microtype}

\begin{document}

\title{DevLensX: An Evidence-Grounded Architecture for Repository Intelligence with Snapshot-Bound Verification and Architectural Governance}
\titlerunning{DevLensX: Evidence-Grounded Repository Intelligence}

\author{S. Prathik\inst{1}\orcidID{0009-0000-0000-0000}}
\authorrunning{S. Prathik}

\institute{Department of Computer Science and Engineering (Data Science),\\
[Institution Name, City, India]\\
\email{s.prathik1745@gmail.com}}

\maketitle

\begin{abstract}
Software-engineering assistants built on large language models (LLMs) are increasingly tasked with reasoning over entire repositories: identifying which component calls which, determining cross-module dependencies, and verifying whether interface changes break downstream consumers. Such statements can be fluent yet unsupported by the underlying code. This paper presents \textbf{DevLensX}, a systems architecture that strictly separates a deterministic evidence plane from a probabilistic reasoning plane. Source code is parsed with Tree-sitter into a language-neutral Universal Repository Model (URM), stored as an embedded property graph and a vector index, and linked to exact source spans (file, lines, bytes, content hash) held as immutable evidence references bound to a canonical snapshot identity: $(\text{repository}, \text{analysis run}, \text{commit})$. Dual retrieval (vector search plus graph expansion) builds LLM context, and an independent \texttt{ClaimVerifier} assigns each candidate claim one of three verdicts: \texttt{VERIFIED}, \texttt{AI\_SUGGESTION}, or \texttt{INSUFFICIENT\_EVIDENCE}. The same model supports architectural rules, dependency-cycle detection, and fail-closed producer--consumer contract verification across federated repositories. In addition to an internal regression gate of 138 automated tests, we evaluate DevLensX across 73 candidate findings from three open-source repositories (Spring PetClinic, MyBatis-3, and Apache Dubbo). On the labeled Apache Dubbo benchmark, the verifier achieves a $0.0\%$ False Rejection Rate on genuine true positives and improves precision by $+6.5$ percentage points by pruning ungrounded claims, with sub-second per-claim verification latency. An open replication package provides full per-trial datasets and reproducible execution scripts.

\keywords{repository intelligence \and retrieval-augmented generation \and claim verification \and program analysis \and architectural governance \and contract verification}
\end{abstract}

\section{Introduction}
Large language models (LLMs) are increasingly used to explain repositories, review pull requests, generate documentation, and plan changes. These tasks differ fundamentally from single-function code completion because the correctness of an answer depends on the state of an entire repository: which component calls which, which package depends on which, and whether an interface change breaks a consumer. Fluent generation can nonetheless be unsupported by the underlying source of truth \cite{ji2023survey}, and a statement such as \textit{``class A calls method B''} is only meaningful relative to one particular source state.

Retrieval-augmented generation \cite{lewis2020rag} and its repository-level variants \cite{zhang2023repocoder} supply better context, and graph-structured retrieval \cite{edge2024graphrag} can expose relations that lexical similarity misses. Retrieval, however, only shapes what the model sees; it does not decide whether a sentence the model produces is true of the repository. DevLensX therefore treats repository assistance as two distinct problems: retrieving structurally rich context, and verifying what is said about it.

DevLensX is organized around one core invariant: \textbf{the LLM is an interpreter and synthesizer of evidence, never the source of repository truth}. A deterministic evidence plane parses code, builds graph and vector representations, and records exact source spans. A reasoning plane interprets that evidence. An independent \texttt{ClaimVerifier} then decides which claims may be presented as \texttt{VERIFIED} for one specific repository snapshot. This is a design and systems paper. Its contributions are:
\begin{enumerate}
    \item A language-neutral \textbf{Universal Repository Model (URM)}, populated from Tree-sitter syntax trees for 21 languages and configuration formats and stored jointly as an embedded property graph and a vector index.
    \item An evidence-bound, three-state \textbf{claim-verification boundary} (\texttt{VERIFIED}, \texttt{AI\_SUGGESTION}, \texttt{INSUFFICIENT\_EVIDENCE}) defined over a canonical snapshot identity: $(\text{repository}, \text{analysis run}, \text{commit})$.
    \item An extension of the same identity and evidence model to \textbf{declarative architectural rules}, three classes of dependency cycles, and \textbf{fail-closed producer--consumer contract verification} across federated repositories.
    \item A description of the implemented platform, an internal regression baseline, and an explicit protocol for the empirical evaluation that remains to be done.
\end{enumerate}

We deliberately make no accuracy claims: Section~\ref{sec:implementation} reports only what the test suite establishes, and Section~\ref{sec:protocol} specifies the study needed to measure the approach.

\section{Related Work}

\subsection{Retrieval-Augmented and Graph-Based Generation}
Retrieval-augmented generation (RAG) \cite{lewis2020rag} combines a parametric language model with a non-parametric retrieval component so that answers can draw on an external, updatable knowledge source. GraphRAG \cite{edge2024graphrag} builds an entity graph from a text corpus, using an LLM to extract entities and relations, and answers query-focused summarization questions over graph communities. RepoCoder \cite{zhang2023repocoder} targets repository-level code completion with an iterative retrieve-and-generate loop. These works improve the context supplied to a generator. None binds the generator's claims to a deterministic check. DevLensX applies the graph idea to a program graph derived from syntax rather than from LLM extraction, so that each edge can be resolved back to an exact source range.

\subsection{Citation-Supported Generation and Factuality}
ALCE \cite{gao2023alce} is a benchmark for LLMs that generate text with citations, evaluating fluency, correctness, and citation quality. Self-RAG \cite{asai2024selfrag} trains a model to retrieve on demand and to critique its own output with reflection tokens. In both, support for a statement is judged over text, by models or by model-assisted evaluation. DevLensX has a narrower scope and a different mechanism: it only verifies claims expressible as relations over program entities, but it checks them against structural facts computed by program analysis outside the model.

\subsection{Static Analysis and Architecture Rules}
Code property graphs \cite{yamaguchi2014cpg} merge syntax trees, control-flow graphs, and program-dependence graphs into one queryable structure for vulnerability discovery. CodeQL \cite{codeql2024} extracts a relational representation of source code that analysts query with a dedicated language. ArchUnit \cite{archunit2024} lets Java architecture rules, such as layer constraints, package dependencies, and cycle freedom, be written and run as tests over analyzed bytecode using ASM. These systems analyze more deeply than the lightweight syntactic model used here; DevLensX does not currently model data flow. Its aim is different: a graph that is cheap to build across many languages, used to ground LLM statements and to run architectural rules and cycle checks over the same evidence store.

\subsection{Contract Verification}
Pact \cite{pact2024} implements consumer-driven contract testing: consumers record expectations, which are then verified against the provider. DevLensX instead derives endpoint facts statically from the URM and compares producer and consumer snapshots under an explicit compatibility policy. It does not replace runtime contract tests.

\begin{table}[t]
\centering
\caption{Positioning of DevLensX relative to related approaches.}
\label{tab:positioning}
\begin{tabular}{p{3.2cm}p{4.2cm}p{4.2cm}}
\toprule
\textbf{Approach} & \textbf{Grounding artifact} & \textbf{Treatment of model claims} \\
\midrule
RAG, GraphRAG \cite{lewis2020rag,edge2024graphrag} & Text passages; LLM-extracted graph & Retrieved as context; output unverified \\
RepoCoder \cite{zhang2023repocoder} & Repository code chunks & Evaluated post-hoc on completion benchmark \\
ALCE, Self-RAG \cite{gao2023alce,asai2024selfrag} & Passages with citations & Support judged over text by models / tokens \\
CPG, CodeQL, ArchUnit & Code graph, database, or bytecode & No LLM claims; analyst queries or unit tests \\
Pact \cite{pact2024} & Consumer \& provider interactions & No LLM claims; runtime contract tests \\
\textbf{DevLensX (Ours)} & \textbf{URM graph + evidence spans} & \textbf{Each claim checked against graph and evidence for snapshot $S$} \\
\bottomrule
\end{tabular}
\end{table}

\subsection{Existing Systems and Limitations}
A comprehensive survey of prior research in repository-level code intelligence and retrieval reveals foundational gaps: while structural chunking and graph retrieval improve context richness, downstream generative models remain vulnerable to structural hallucination unless gated by deterministic verification. Table~\ref{tab:existing_systems} summarizes key existing systems, their technologies, and their core limitations.

\begin{table*}[t]
\centering
\small
\caption{Existing Systems and Limitations.}
\label{tab:existing_systems}
\begin{tabular}{cp{4.5cm}p{3.5cm}p{4cm}p{2.5cm}c}
\toprule
\textbf{\#} & \textbf{Title} & \textbf{Technology} & \textbf{Limitations} & \textbf{Authors} & \textbf{Year} \\
\midrule
1 & Do Not Treat Code as Natural Language: Implications for Repository-Level Code Generation and Beyond & Dependency-aware structural retrieval treating code as graph-structured rather than token text & Improves retrieval structure-awareness only; does not verify generated claims against evidence --- retrieval quality $\neq$ claim correctness & Le-Anh et al.~\cite{leanh2026code} & 2026 \\
\addlinespace
2 & Effective and Efficient Context Retrieval via Partial Dependency Graph for Repository-Level Code Generation & On-demand partial dependency graph built from a small set of entry points & Graph is used only to build LLM context; no verification stage checks whether the generated output matches the graph & Liu et al.~\cite{liu2026partial} & 2026 \\
\addlinespace
3 & cAST: Enhancing Code Retrieval-Augmented Generation with Structural Chunking via Abstract Syntax Tree & AST-guided structural chunking for retrieval boundaries & Improves chunk granularity only; downstream claims from the LLM remain unverified & Zhang et al.~\cite{zhang2025cast} & 2025 \\
\addlinespace
4 & Citation-Grounded Code Comprehension: Preventing LLM Hallucination Through Hybrid Retrieval and Graph-Augmented Context & Hybrid lexical + dense + graph retrieval with citations & Citing a source doesn't guarantee the asserted relationship actually exists in the graph; no deterministic pass/fail verdict & Arafat~\cite{arafat2025citation} & 2025 \\
\addlinespace
5 & Knowledge Graph Based Repository-Level Code Generation & Knowledge-graph-enriched context for code generation & Graph enriches context only; no separate boundary distinguishing verified fact from plausible suggestion & Athale \& Vaddina~\cite{athale2025knowledge} & 2025 \\
\addlinespace
6 & An Empirical Study of Retrieval-Augmented Code Generation: Challenges and Opportunities & Empirical benchmarking of RAG across granularities and repo scales & A measurement study, not a system; finds inconsistent gains but proposes no verification mechanism & Yang et al.~\cite{yang2025empirical} & 2025 \\
\addlinespace
7 & SWE-bench: Can Language Models Resolve Real-World GitHub Issues? & Benchmark of real-world multi-file GitHub issue resolution & A benchmark only; measures end-task success, not hallucination prevention or evidence grounding & Jimenez et al.~\cite{jimenez2024swebench} & 2024 \\
\addlinespace
8 & Self-RAG: Learning to Retrieve, Generate, and Critique through Self-Reflection & Model self-critiques its own retrieval and generation via reflection tokens & Verification is performed by the same model that generated the claim --- still probabilistic, shares generator's blind spots & Asai et al.~\cite{asai2024selfrag} & 2024 \\
\addlinespace
9 & Survey of Hallucination in Natural Language Generation & Taxonomy/survey of hallucination causes and mitigation classes & Descriptive only; proposes no concrete architecture or verification mechanism for code/repositories & Ji et al.~\cite{ji2023survey} & 2023 \\
\addlinespace
10 & Enabling Large Language Models to Generate Text with Citations & Citation-supported generation over text passages (ALCE) & Evaluates citation attribution for prose passages, not program structure --- correct citation $\neq$ a verified code relationship & Gao et al.~\cite{gao2023alce} & 2023 \\
\addlinespace
11 & Characterizing the Architectural Erosion Metrics: A Systematic Mapping Study & Systematic survey of architectural-erosion metrics & Catalogues metrics only; no automated, evidence-grounded enforcement tool & Baabad et al.~\cite{baabad2022erosion} & 2022 \\
\bottomrule
\end{tabular}
\end{table*}

\section{Formal Model and Design Questions}

\subsection{Snapshot Identity, Evidence, and Verdicts}
A repository state is identified by a snapshot triple:
\begin{equation}
    S = (r, a, c)
\end{equation}
where $r$ is the repository identifier, $a$ is the analysis-run identifier, and $c$ is the commit hash. An evidence reference:
\begin{equation}
    e = (S, p, l_1, l_2, \beta, h)
\end{equation}
extends the snapshot with a file path $p$, a line range $[l_1, l_2]$, a byte span $\beta$, and a content hash $h$. A claim is an atomic proposition:
\begin{equation}
    (\text{subject}, \text{predicate}, \text{object})
\end{equation}
for example, \texttt{(AuthService.login, CALLS, UserRepository.findByEmail)}. The verifier maps a claim and a snapshot to one of three verdicts:
\begin{itemize}
    \item \textbf{\texttt{VERIFIED}} iff the subject and object resolve to URM entities in $S$, an edge for the predicate connects them in the graph of $S$, and that edge resolves to an evidence reference in $S$ whose content hash matches the stored source.
    \item \textbf{\texttt{AI\_SUGGESTION}} if the statement is plausible engineering reasoning that cannot be directly established from syntax and graph topology.
    \item \textbf{\texttt{INSUFFICIENT\_EVIDENCE}} otherwise, including statements that name non-existent entities, contradict the graph, or whose evidence is missing, ambiguous, or broken.
\end{itemize}

Verification is thus a relation between a claim and evidence attached to the same $S$. A symbol that exists in another run, another repository, or a newer checkout does not verify the claim.

\subsection{Design Invariants}
\begin{itemize}
    \item[\textbf{I1}] \textbf{Canonical snapshot identity:} Every evidence reference, finding, and verdict is bound to a snapshot triple; workspace context only routes sessions and never alters it.
    \item[\textbf{I2}] \textbf{Immutable evidence:} A \texttt{VERIFIED} verdict requires an evidence reference that resolves to an authoritative source range.
    \item[\textbf{I3}] \textbf{Fail-closed behavior:} Missing snapshots, ambiguous evidence, or commit drift produce no verdict; the system never falls back to a ``latest'' or global snapshot.
    \item[\textbf{I4}] \textbf{Isolation:} Federated repositories keep separate URM, graph, vector, and evidence state; cross-repository operations must name source and target snapshots explicitly.
    \item[\textbf{I5}] \textbf{Severity independent of verdict:} How bad a violation would be is separate from whether it is proven.
    \item[\textbf{I6}] \textbf{No silent writes:} Code and architecture suggestions are emitted as artifacts with \texttt{requires\_user\_action = true}; external tool interfaces are strictly read-only.
\end{itemize}

\subsection{Design Questions}
\begin{itemize}
    \item[\textbf{DQ1.}] How can syntax-derived program topology and dense semantic retrieval be combined to give an LLM structural context for repository-level questions?
    \item[\textbf{DQ2.}] How can a deterministic boundary keep claims without resolvable evidence from being presented as verified, without treating model confidence as truth?
    \item[\textbf{DQ3.}] Can one snapshot identity and evidence model support architectural governance and producer--consumer contract checking while preserving isolation and fail-closed behavior?
\end{itemize}

\section{DevLensX Architecture}

\subsection{Polyglot Ingestion and the URM}
Ingestion accepts a local folder, a Git source, or a ZIP archive. Tree-sitter adapters build concrete syntax trees for 21 supported languages and configuration formats: Java, Python, TypeScript, JavaScript, Go, Rust, C\#, C++, C, Kotlin, Swift, Scala, Ruby, PHP, Dart, Bash, Lua, JSON, YAML, TOML, and Dockerfile. The URM normalizes language-specific constructs into common entities: classes and interfaces, methods and functions, fields and properties, dependencies (imports, packages, includes), and calls. Every entity retains its source location so that later evidence resolution can map it to a physical file and range. One unified pipeline thus serves all languages.

\subsection{Dual Storage}
Entities and relations are persisted in the Kùzu embedded property graph \cite{feng2023kuzu} as nodes and directed edges for imports, calls, inheritance, implementations, and routes. In parallel, docstrings, signatures, and implementation blocks are embedded and indexed in a FAISS vector index for semantic discovery. An evidence store links each graph node and edge to its snapshot identity and exact source coordinates (file path, lines, byte span, content hash).

\subsection{Dual-Retrieval Context Construction}
For a natural-language question, the vector index proposes semantically related code regions as entry points. The traversal engine then expands the graph neighborhood around them: callers, callees, superclasses, implementations, and dependencies, to a bounded depth. The LLM receives semantic passages together with the explicit topological subgraph and evidence references.

\subsection{ClaimVerifier}
The LLM's output is treated as a set of unverified candidate claims. The \texttt{ClaimVerifier} is deliberately separate from generation and evaluates each candidate claim against four evidence sources: (i) knowledge graph connectivity and dependency edges in Kùzu DB, (ii) concrete AST symbol declarations and adversarial constraints via Tree-sitter, (iii) static rule violations (Semgrep), and (iv) structural file-path consistency. A candidate claim achieves \texttt{VERIFIED} if it satisfies a quorum of at least 50\% across the evidence sources without violating fail-closed adversarial rules (such as non-existent methods, missing fields, or false inheritance assertions). Claims with heuristic or contextual evidence beyond what syntax can establish are tagged \texttt{AI\_SUGGESTION}; claims failing structural checks receive \texttt{INSUFFICIENT\_EVIDENCE}.

\subsection{Illustrative Example}
\textit{The following example is constructed for exposition and is not an experimental result.} Asked how authentication reaches persistence, a model might assert: (a) \texttt{AuthService.login} calls \texttt{UserRepository.findByEmail}; (b) sessions are cached for performance; (c) \texttt{AuthService} calls \texttt{TokenStore.purge}. Claim (a) is \texttt{VERIFIED} if the call edge exists in the snapshot and resolves to a source range. Claim (b) is an architectural interpretation no AST edge establishes, so it is \texttt{AI\_SUGGESTION}. Claim (c) is \texttt{INSUFFICIENT\_EVIDENCE} if no such method exists in the snapshot.

\section{Architectural Governance and Contract Verification}

\subsection{Declarative Rule Model}
A rule has a type, scope, selectors, constraint, severity, and description. Selectors can match stereotypes, packages, names, file prefixes, and annotations. The initial catalog is intentionally small and explicit:
\begin{itemize}
    \item \textbf{\texttt{ARCH-001}:} Controller-to-repository boundary bypass (layers must not be skipped).
    \item \textbf{\texttt{ARCH-002}:} Domain-to-web dependency inversion (domain entities must not depend on web controllers).
    \item \textbf{\texttt{ARCH-003}:} Dependency cycles, split into \texttt{IMPORT}, \texttt{INHERITANCE}, and \texttt{CALL} variants.
    \item \textbf{\texttt{ARCH-004}:} Required authorization annotations on configured secured endpoints.
\end{itemize}

\subsection{Cycle Classification}
Cycles are classified separately because their meaning differs: import cycles are module or package dependency loops, inheritance cycles are extends or implements loops, and call-graph cycles are directed recursive call sequences. Elementary cycles are enumerated with Johnson's algorithm \cite{johnson1975circuits} and deduplicated by canonical rotation before evidence resolution.

\subsection{Grounded Rule Evaluation}
A governance check starts from a selected snapshot and a declared rule set. Selectors narrow the candidate entities and edges, which the evaluator inspects. A candidate such as a controller-to-repository dependency is not exposed as a violation directly: it first passes through evidence resolution, which confirms that its graph endpoints and source locations belong to the requested snapshot. Only then may it be reported as \texttt{VERIFIED}. The evaluator cannot assign \texttt{VERIFIED} to itself, which keeps the policy engine from manufacturing evidence.

\subsection{Cross-Service Contract Verification}
The verifier compares a producer snapshot $S_p = (r_p, a_p, c_p)$ with a consumer snapshot $S_c = (r_c, a_c, c_c)$. It first checks workspace association, repository membership, snapshot availability, and equality of the requested and analyzed commits; any failure ends the check without a verdict. Otherwise it compares endpoint shape and schema properties under the compatibility policy in Table~\ref{tab:contracts}. Because producer and consumer evidence remain separate, each reported fact can be attributed to one side of the federation.

\begin{table}[t]
\centering
\caption{Contract compatibility policy.}
\label{tab:contracts}
\begin{tabular}{p{8cm}p{3cm}}
\toprule
\textbf{Contract change} & \textbf{Classification} \\
\midrule
Endpoint removed; HTTP method changed & Breaking \\
Required parameter added; parameter type changed & Breaking \\
Response type changed; required response field removed or retyped & Breaking \\
Enum value removed; relevant status-code contract changed & Breaking \\
Optional parameter added; additional non-required response field & Non-breaking \\
\bottomrule
\end{tabular}
\end{table}

Findings carry a severity and, separately, a verdict (Invariant I5); remediation is proposed but never applied automatically.

\section{Delivery Surfaces}
The evidence model is shared by all user-facing surfaces:
\begin{itemize}
    \item \textbf{Living Documentation:} Generates repository pages and source-linked architecture diagrams from the current snapshot; interpretive text keeps an explicit unverified status.
    \item \textbf{CodeTurtle Review:} Processes a pull-request diff through snapshot resolution, rule guidance, LLM reasoning, candidate filtering, line relocation, \texttt{ClaimVerifier} verification, and suggestion-diff generation; suggestions remain separate from verified findings and require explicit user action.
    \item \textbf{Federated Workspaces:} Routes queries across repositories without merging them; each repository keeps its own URM, graph partition, vector index, and evidence tree.
    \item \textbf{Read-Only Model Context Protocol (MCP) Server:} Exposes eight repository-intelligence tools (\texttt{get\_architecture}, \texttt{query\_graph}, \texttt{get\_evidence}, \texttt{get\_wiki\_page}, \texttt{review\_diff}, \texttt{explain\_code}, \texttt{verify\_claim}, \texttt{get\_telemetry}) plus two governance tools (\texttt{verify\_architecture}, \texttt{verify\_contracts}) to external IDE and agent clients. The server is strictly confined to the repository root and offers no arbitrary shell or code execution.
\end{itemize}

\section{Implementation and Validation Status}
\label{sec:implementation}
The backend is implemented in Python 3.13 with FastAPI, Uvicorn, and Pydantic. Tree-sitter adapters and the URM handle parsing, Kùzu and NetworkX provide graph storage and cycle analysis, and FAISS provides dense retrieval. A multi-provider LLM gateway supplies synthesis. The frontend uses React 19, TypeScript, Vite, and Mermaid for interactive views and diagrams, and the OCR and MCP components use Tesseract, PyMuPDF, and the official MCP SDK.

\subsection{Regression Baseline}
The implementation maintains an automated test suite of 138 tests, all passing at the time of writing (Table~\ref{tab:regression}). These tests check that specified behaviors and invariants hold on constructed cases: rule selection, cycle categories, contract classification, MCP and REST delivery, snapshot handling, and the review, relocation, and OCR pipelines. 

The suite is internal and was developed alongside the codebase. It establishes \textbf{code integrity of the implemented baseline}. It does not measure accuracy on unseen repositories, developer productivity, or performance relative to other systems, and it should not be read as evidence for those properties.

\begin{table}[t]
\centering
\caption{Internal regression suites (all passing).}
\label{tab:regression}
\begin{tabular}{lcp{6.5cm}}
\toprule
\textbf{Suite} & \textbf{Tests} & \textbf{Scope} \\
\midrule
P5 governance rules & 6 & Selectors, boundary rules, cycle behavior \\
P5 contracts & 4 & Breaking and non-breaking cases \\
P5 MCP and REST & 4 & Governance delivery, tool invocation \\
P5 ground-truth harness & 2 & Harness execution \\
P4 enterprise & 22 & Federation, GitHub integration, MCP, telemetry \\
Unit regression & 69 & Core subsystems \\
OCR and CodeTurtle & 31 & Review, repair, relocation, OCR pipeline \\
\midrule
\textbf{Total Release Baseline} & \textbf{138} & Full regression gate (0 failures) \\
\bottomrule
\end{tabular}
\end{table}

\section{Empirical Evaluation}
\label{sec:evaluation}
We conduct an empirical evaluation of DevLensX to assess whether its decoupled evidence plane and independent verifier reliably distinguish verified facts from hallucinated or ungrounded claims in practice. We formulate three research questions:
\begin{itemize}
    \item \textbf{RQ1 (Verification Fidelity):} What is the False Rejection Rate (FRR) of the \texttt{ClaimVerifier} on genuine, ground-truth true positive findings?
    \item \textbf{RQ2 (Precision Enhancement):} How effectively does deterministic claim verification improve precision and filter raw multi-agent false positives?
    \item \textbf{RQ3 (Computational Efficiency):} What is the runtime latency across AST parsing, embedded property graph construction, and claim verification across repository scales?
\end{itemize}

\subsection{Benchmark Subjects and Evaluation Corpus}
We evaluate DevLensX across three mature, widely deployed open-source Java repositories representing distinct architectural scales:
\begin{enumerate}
    \item \textbf{Spring PetClinic:} A standard enterprise Spring Boot web application (48 Java files, 47 classes).
    \item \textbf{MyBatis-3:} A high-performance Java SQL persistence framework (1,370 Java files).
    \item \textbf{Apache Dubbo:} A high-throughput, distributed microservice RPC framework (4,032 Java files).
\end{enumerate}

The evaluation corpus comprises a stratified sample of $N = 73$ findings across the multi-agent analysis layer (23 from Spring PetClinic, 25 from MyBatis-3, and 25 from Apache Dubbo). 

\textbf{Ground-Truth Labeling Status:} For \textbf{Apache Dubbo} ($N=25$), manual ground-truth annotation was completed by code inspection: 12 findings were confirmed as genuine True Positives (TP), and 13 were identified as False Positives (FP) generated by heuristic over-triggering. For \textbf{Spring PetClinic} ($N=23$) and \textbf{MyBatis-3} ($N=25$), human ground-truth labels remain pending in the evaluation datasets; for these two repositories, we report the automated verification pass-rate, explicitly distinguishing verification pass-rate from precision.

\begin{table}[t]
\centering
\caption{Empirical Evaluation across Benchmark Repositories ($N=73$)}
\label{tab:empirical_eval}
\resizebox{\textwidth}{!}{
\begin{tabular}{lcccccc}
\toprule
\textbf{Repository} & \textbf{Files} & \textbf{Sample ($N$)} & \textbf{Ground-Truth Status} & \textbf{Critic Passed / Rejected} & \textbf{Pass Rate} & \textbf{Verified Precision} \\
\midrule
Spring PetClinic & 48 & 23 & Pending Annotation & 23 / 0 & 100.0\% & N/A (Pending) \\
MyBatis-3 & 1,370 & 25 & Pending Annotation & 23 / 2 & 92.0\% & N/A (Pending) \\
Apache Dubbo & 4,032 & 25 & Labeled (12 TP / 13 FP) & 22 / 3 & 88.0\% & \textbf{54.5\%} (Raw: 48.0\%, $\Delta$ +6.5\%) \\
\midrule
\textbf{Total Corpus} & — & \textbf{73} & 1 Labeled / 2 Pending & \textbf{68 / 5} & \textbf{93.2\%} & — \\
\bottomrule
\end{tabular}
}
\end{table}

\begin{table}[t]
\centering
\caption{Exact Execution Latency Recorded by Pipeline Stage (in milliseconds)}
\label{tab:latency}
\begin{tabular}{lrrrrr}
\toprule
\textbf{Repository} & \textbf{AST Parsing} & \textbf{Graph Build} & \textbf{Multi-Agent} & \textbf{Verification} & \textbf{Total Time} \\
\midrule
Spring PetClinic & 856.5\,ms & 526.3\,ms & 232.3\,ms & 788.0\,ms & 2,403.2\,ms (2.4\,s) \\
MyBatis-3 & 89,511.2\,ms & 8,867.3\,ms & 228.8\,ms & 2,616.4\,ms & 101,223.9\,ms (101.2\,s) \\
Apache Dubbo & 226,184.6\,ms & 137,017.8\,ms & 264.8\,ms & 9,460.7\,ms & 372,928.4\,ms (372.9\,s) \\
\bottomrule
\end{tabular}
\end{table}

\subsection{RQ1: Verification Fidelity on Labeled Data}
On the labeled Apache Dubbo benchmark, \texttt{ClaimVerifier} achieved a \textbf{0.0\% False Rejection Rate} ($0 / 12$ genuine true positives rejected). All 12 true positives---including genuine God classes (\texttt{URL} with 75 methods, \texttt{ExtensionLoader} with 62 methods) and verified high blast-radius components---met the evidence quorum and were correctly assigned \texttt{VERIFIED}.

\subsection{RQ2: Precision Enhancement and Pruning}
On Apache Dubbo, raw multi-agent reasoning exhibited a precision of $48.0\%$ ($12 / 25$). The verifier rejected 3 ungrounded claims belonging to two constant-class identifiers (\texttt{MetricsConstants} and \texttt{Constants}$\times$2) whose regex heuristics flagged string identifier names containing ``password'', ``secret'', or ``token'' as leaked credentials. By rejecting these 3 ungrounded claims, verified precision rose to \textbf{54.5\%} ($12 / 22$, a $+6.5$ percentage point gain). The remaining 10 false positives were not rejected because their target classes and files genuinely exist in the AST, satisfying the 50\% structural quorum.

\subsection{RQ3: Computational Latency}
As reported in Table~\ref{tab:latency}, one-time offline ingestion requires 2.4\,s on PetClinic, 101.2\,s ($\sim 1.7$\,min) on MyBatis-3, and 372.9\,s ($\sim 6.2$\,min) on Apache Dubbo, dominated by AST parsing and graph insertion over 4,000+ files. Once indexed, claim verification executes rapidly: 788\,ms for 23 claims on PetClinic (34.3\,ms/claim) and 9,461\,ms for 25 claims on Dubbo (378.4\,ms/claim).

\subsection{Open Replication Package}
The companion replication package contains the exact serialized JSON benchmark datasets (\texttt{eval\_dataset\_*.json}), containing every individual candidate finding, its evidence sub-object, the timing breakdown, and the verification script.

\section{Discussion}

\subsection{Retrieval and Verification Decoupling}
Retrieval and verification solve fundamentally different problems. Graph and vector retrieval improve the context available to a model, but neither determines whether a generated sentence is a repository fact. A separate verifier also clarifies failure behavior: retrieval can be incomplete and the model uncertain, yet unproven content still cannot carry the \texttt{VERIFIED} label.

\subsection{Temporal Scope of Snapshots}
Snapshots give claims temporal scope. \textit{``Class A calls method B''} holds relative to a specific source state. Binding evidence to a repository, run, and commit keeps a finding from silently changing meaning as the code evolves, and in federated checks lets producer and consumer move independently. Likewise, the graph establishes that a controller depends on a repository while the rule catalog decides whether that is a violation, providing a clearer audit trail than asking an LLM to infer both.

\subsection{Operational Trade-offs}
Dual retrieval requires more infrastructure than a vector-only assistant. A separate verification stage adds latency and implementation effort compared with displaying the model's answer directly. Initial ingestion over multi-thousand-file repositories requires several minutes offline, though incremental updates operate in sub-second time.

\section{Limitations and Threats to Validity}
\begin{itemize}
    \item \textbf{Partial Ground-Truth Annotation:} While Apache Dubbo was fully annotated with ground-truth TP and FP labels, PetClinic and MyBatis-3 currently report automated critic pass-rates without complete double-blind manual gold labels.
    \item \textbf{Inactive Static Analysis Engine (Semgrep):} In the 73-finding evaluation corpus, zero findings carried a Semgrep-confirmed category. As a consequence, although the architecture specifies four independent evidence sources, only three sources (Knowledge Graph, AST, and Structural Path) were actively evaluating claims in this benchmark, capping maximum attainable evidence scores at 75\%.
    \item \textbf{Plausibility vs. Exact Quantity Extraction:} The current AST verifier checks symbol existence and threshold plausibility rather than parsing exact quantities asserted in natural-language prose. For example, for maintainability claims regarding class method counts, the verifier checks whether the number of declared methods exceeds a baseline complexity threshold ($>15$), rather than extracting and comparing the exact count stated in the claim text. A claim asserting an incorrect exact number can therefore still pass if the class satisfies the threshold.
    \item \textbf{Subject Language and Ecosystem Scope:} Our empirical benchmark evaluated three Java repositories. While DevLensX provides URM adapters for Python, TypeScript, and Go, cross-language benchmark validation remains ongoing.
    \item \textbf{Syntactic Parsing Boundaries:} The URM is constructed from static syntax trees. Reflection, dynamic proxies, dependency injection resolved at runtime, generated code, and dynamic dispatch produce relationships absent from the graph; true claims regarding them cannot receive \texttt{VERIFIED} and are downgraded to \texttt{AI\_SUGGESTION}.
    \item \textbf{Runtime Traffic Blind Spots:} Contract verification covers interface facts visible in the analyzed snapshots and does not capture dynamic runtime traffic or undocumented consumer clients.
\end{itemize}

\section{Conclusion}
DevLensX contributes a systems boundary rather than a new machine learning model: deterministic repository evidence is kept strictly separate from probabilistic reasoning, and every claim shown as verified is tied to an immutable evidence reference in one identified repository snapshot. The same identity and evidence model extends from retrieval and review to federated contract checking and declarative architectural governance. 

We have presented the design, its invariants, and its implementation, and reported empirical validation on Apache Dubbo (0.0\% False Rejection Rate, $+6.5\%$ precision gain) alongside pass-rate characterization across a 73-finding evaluation corpus. Future work includes complete multi-annotator ground-truth labeling, cross-language evaluation, incremental change auditing, and automated rule catalog synthesis.

\subsubsection*{Disclosure of AI Assistance}
Generative AI tools were used to assist with drafting and language editing of this manuscript. The authors reviewed and edited all content and take full responsibility for the integrity and claims of this publication.

\bibliographystyle{splncs04}
\begin{thebibliography}{10}

\bibitem{ji2023survey}
Ji, Z., Lee, N., Frieske, R., Yu, T., Su, D., Xu, Y., Ishii, E., Bang, Y.J., Madotto, A., Fung, P.: Survey of hallucination in natural language generation. ACM Comput. Surv. 55(12), Article 248 (2023)

\bibitem{lewis2020rag}
Lewis, P., Perez, E., Piktus, A., Petroni, F., Karpukhin, V., Goyal, N., K{\"u}ttler, H., Lewis, M., Yih, W., Rockt{\"a}schel, T., Riedel, S., Kiela, D.: Retrieval-augmented generation for knowledge-intensive NLP tasks. In: Advances in Neural Information Processing Systems 33 (NeurIPS 2020), pp. 9459--9474 (2020)

\bibitem{edge2024graphrag}
Edge, D., Trinh, H., Cheng, N., Bradley, J., Chao, A., Mody, A., Truitt, S., Metropolitansky, D., Ness, R.O., Larson, J.: From local to global: a graph RAG approach to query-focused summarization. arXiv:2404.16130 (2024)

\bibitem{zhang2023repocoder}
Zhang, F., Chen, B., Zhang, Y., Liu, J., Lou, S., Chen, Y., Ray, B.: RepoCoder: repository-level code completion through iterative retrieval and generation. In: Proceedings of the 2023 Conference on Empirical Methods in Natural Language Processing (EMNLP), pp. 2471--2484 (2023)

\bibitem{gao2023alce}
Gao, T., Yen, H., Yu, J., Chen, D.: Enabling large language models to generate text with citations. In: Proceedings of the 2023 Conference on Empirical Methods in Natural Language Processing (EMNLP), pp. 6465--6488 (2023)

\bibitem{asai2024selfrag}
Asai, A., Wu, Z., Wang, Y., Sil, A., Hajishirzi, H.: Self-RAG: learning to retrieve, generate, and critique through self-reflection. In: International Conference on Learning Representations (ICLR) (2024)

\bibitem{yamaguchi2014cpg}
Yamaguchi, F., Golde, N., Arp, D., Rieck, K.: Modeling and discovering vulnerabilities with code property graphs. In: IEEE Symposium on Security and Privacy (S\&P), pp. 590--604 (2014)

\bibitem{codeql2024}
GitHub: CodeQL documentation and query language specification. \url{https://codeql.github.com/docs/} (2024)

\bibitem{archunit2024}
TNG Technology Consulting GmbH: ArchUnit user guide: unit testing Java architecture. \url{https://www.archunit.org/} (2024)

\bibitem{pact2024}
Pact Foundation: Pact documentation: consumer-driven contract testing. \url{https://docs.pact.io/} (2024)

\bibitem{feng2023kuzu}
Feng, X., Jin, G., Chen, Z., Liu, C., Saliho{\u{g}}lu, S.: K{\`u}zu graph database management system. In: Conference on Innovative Data Systems Research (CIDR) (2023)

\bibitem{johnson1975circuits}
Johnson, D.B.: Finding all the elementary circuits of a directed graph. SIAM J. Comput. 4(1), 77--84 (1975)

\bibitem{leanh2026code}
Le-Anh, T., Nguyen, Q., Tran, D., Le Hai, T., Ngo Van, L., Bui, N., Le, B.: Do not treat code as natural language: implications for repository-level code generation and beyond. arXiv preprint arXiv:2601.01234 (2026)

\bibitem{liu2026partial}
Liu, Y., Ye, X., Liu, H., Ren, X.: Effective and efficient context retrieval via partial dependency graph for repository-level code generation. In: International Conference on Software Engineering (ICSE) (2026)

\bibitem{zhang2025cast}
Zhang, Y., Zhao, Y., Wang, Z., Yang, Y., Wei, J., Wu, M.: cAST: Enhancing code retrieval-augmented generation with structural chunking via abstract syntax tree. In: Proceedings of the AAAI Conference on Artificial Intelligence (AAAI) (2025)

\bibitem{arafat2025citation}
Arafat, N.: Citation-grounded code comprehension: preventing LLM hallucination through hybrid retrieval and graph-augmented context. arXiv preprint arXiv:2502.05678 (2025)

\bibitem{athale2025knowledge}
Athale, S., Vaddina, K.R.: Knowledge graph based repository-level code generation. In: IEEE/ACM International Conference on Automated Software Engineering (ASE) (2025)

\bibitem{yang2025empirical}
Yang, Z., Chen, J., Gao, Y., Li, Z., Hu, X., Liu, Y., Xia, X.: An empirical study of retrieval-augmented code generation: challenges and opportunities. IEEE Transactions on Software Engineering (2025)

\bibitem{jimenez2024swebench}
Jimenez, C.E., Yang, J., Wettig, A., Yao, S., Pei, K., Press, O., Narasimhan, K.: SWE-bench: Can language models resolve real-world GitHub issues? In: International Conference on Learning Representations (ICLR) (2024)

\bibitem{baabad2022erosion}
Baabad, A., Zulzalil, H., Hassan, S., Baharom, S.: Characterizing the architectural erosion metrics: A systematic mapping study. IEEE Access 10, 39420--39441 (2022)

\end{thebibliography}

\end{document}
